% ─────────────────────────────────────────────────────────────────────────────
%  LAB 03 CHEATSHEET — Miền tần số & Miền không gian + Mini-Project Pipeline
%  Bùi Quang Chiến — 23001837   |   MAT3562E Computer Vision 2025-2026
% ─────────────────────────────────────────────────────────────────────────────
\documentclass[9pt,a4paper,landscape]{extarticle}

\usepackage[utf8]{inputenc}
\usepackage[T5]{fontenc}
\usepackage[vietnamese,english]{babel}
\usepackage[margin=0.8cm]{geometry}
\usepackage{amsmath,amssymb}
\usepackage{multicol}
\usepackage{booktabs}
\usepackage{xcolor}
\usepackage{listings}
\usepackage{mdframed}
\usepackage{titlesec}
\usepackage{enumitem}
\usepackage{parskip}

\setlength{\parindent}{0pt}
\setlength{\columnsep}{1em}
\setlength{\columnseprule}{0.3pt}
\setlist{nosep, leftmargin=1em}

% ─── Colors ──────────────────────────────────────────────────────────────────
\definecolor{hdr}{RGB}{30,80,150}
\definecolor{subhdr}{RGB}{60,120,200}
\definecolor{codebg}{RGB}{245,245,248}
\definecolor{boxbg}{RGB}{230,242,255}
\definecolor{warnbg}{RGB}{255,243,220}
\definecolor{greenbg}{RGB}{230,255,235}

% ─── Code style ──────────────────────────────────────────────────────────────
\lstset{
  basicstyle=\ttfamily\scriptsize,
  backgroundcolor=\color{codebg},
  frame=none,
  breaklines=true,
  tabsize=2,
  showstringspaces=false,
  keywordstyle=\color{blue!70}\bfseries,
  commentstyle=\color{green!50!black},
  stringstyle=\color{red!60!black},
  aboveskip=1pt, belowskip=1pt,
}

% ─── Note boxes ──────────────────────────────────────────────────────────────
\newmdenv[backgroundcolor=warnbg, linewidth=0.5pt, linecolor=orange!60,
  innerleftmargin=4pt, innerrightmargin=4pt,
  innertopmargin=2pt, innerbottommargin=2pt,
  skipabove=1pt, skipbelow=1pt]{notebox}
\newmdenv[backgroundcolor=greenbg, linewidth=0.5pt, linecolor=green!50!black,
  innerleftmargin=4pt, innerrightmargin=4pt,
  innertopmargin=2pt, innerbottommargin=2pt,
  skipabove=1pt, skipbelow=1pt]{tipbox}
\newmdenv[backgroundcolor=boxbg, linewidth=0.5pt, linecolor=hdr,
  innerleftmargin=4pt, innerrightmargin=4pt,
  innertopmargin=2pt, innerbottommargin=2pt,
  skipabove=1pt, skipbelow=1pt]{infobox}

% ─── Section styles ──────────────────────────────────────────────────────────
\titleformat{\section}{\bfseries\small\color{hdr}}{}{0pt}{}[\vspace{-4pt}\hrule\vspace{2pt}]
\titleformat{\subsection}{\bfseries\scriptsize\color{subhdr}}{}{0pt}{}
\titlespacing{\section}{0pt}{4pt}{2pt}
\titlespacing{\subsection}{0pt}{3pt}{1pt}

% ─────────────────────────────────────────────────────────────────────────────
\begin{document}
\begin{center}
  {\large\bfseries\color{hdr} LAB 03 CHEATSHEET --- Miền Tần Số \& Miền Không Gian + Mini-Project Pipeline} \\[1pt]
  {\small Bui Quang Chien \textbullet{} 23001837 \textbullet{} MAT3562E Computer Vision 2025--2026}
\end{center}
\hrule\vspace{2pt}

\begin{multicols}{3}

% ══════════════════════════════════════════════════════════════════════════════
\section{1. Biến đổi Fourier rời rạc (DFT / FFT)}

\textbf{Công thức DFT 2D:}
\[
  F(u,v) = \sum_{x=0}^{M-1}\sum_{y=0}^{N-1} f(x,y)\, e^{-j2\pi\!\left(\frac{ux}{M}+\frac{vy}{N}\right)}
\]
\textbf{IDFT 2D:}
\[
  f(x,y) = \frac{1}{MN}\sum_{u}\sum_{v} F(u,v)\, e^{j2\pi\!\left(\frac{ux}{M}+\frac{vy}{N}\right)}
\]

\subsection{FFT bằng NumPy}
\begin{lstlisting}[language=Python]
F  = np.fft.fft2(img)        # DFT 2D -> so phuc
Fs = np.fft.fftshift(F)      # DC ra tam
S  = np.log(1 + np.abs(Fs))  # Pho bien do (log scale)
\end{lstlisting}

\begin{notebox}
\textbf{Lưu ý:}
\begin{itemize}
  \item \texttt{fftshift}: dịch DC từ góc ra \textbf{chính giữa}
  \item Dùng \texttt{log(1+|F|)}: DC rất lớn, log chặn tương phản
  \item DC = giá trị trung bình tất cả pixel $\Rightarrow$ điểm sáng nhất ở tâm
  \item Ảnh nhiều texture $\Rightarrow$ phổ lan rộng ra xa tâm hơn
\end{itemize}
\end{notebox}

\textbf{Độ phức tạp:} DFT ngây thơ $O(M^2N^2)$; FFT $O(MN\log MN)$

% ══════════════════════════════════════════════════════════════════════════════
\section{2. Lọc trong miền tần số}

\subsection{Bộ lọc thông thấp lý tưởng (Ideal LPF)}
\begin{lstlisting}[language=Python]
def circular_lpf_mask(shape, r):
    M, N = shape
    crow, ccol = M//2, N//2
    Y, X = np.ogrid[:M, :N]
    dist2 = (Y-crow)**2 + (X-ccol)**2
    return (dist2 <= r*r).astype(np.float64)

Hlp = circular_lpf_mask(img.shape, r=30)
Hhp = 1 - Hlp  # HPF = 1 - LPF
\end{lstlisting}

\subsection{Áp dụng bộ lọc}
\begin{lstlisting}[language=Python]
def apply_freq_filter(image, mask):
    F  = np.fft.fft2(image)
    Fs = np.fft.fftshift(F)
    Gs = Fs * mask                       # Nhan pho voi mat na
    g  = np.fft.ifft2(np.fft.ifftshift(Gs))
    return np.abs(g)
\end{lstlisting}

\begin{center}\small
\begin{tabular}{@{}lll@{}}
\toprule
$r$ & LPF & HPF \\
\midrule
Nhỏ (10) & Rất mờ & Chỉ biên mạnh \\
Vừa (30) & Mờ vừa & Biên rõ \\
Lớn (60) & Gần gốc & Chi tiết nhiều \\
\bottomrule
\end{tabular}
\end{center}

\begin{notebox}
\textbf{Hiệu ứng Ringing (Gibbs):} Ideal LPF cắt ngưỡng đột ngột
$\Rightarrow$ sóng gợn quanh biên. Gaussian LPF mượt hơn, ít ringing hơn.
\end{notebox}

% ══════════════════════════════════════════════════════════════════════════════
\section{3. Tích chập rời rạc 2D}

\textbf{Công thức:}
$(f*h)(x,y) = \displaystyle\sum_{i}\sum_{j} f(x-i,y-j)\cdot h(i,j)$

\textbf{Định lý tích chập:}
$f*h \;\Longleftrightarrow\; F\cdot H$

\begin{lstlisting}[language=Python]
# Tu cai dat (same padding, zero-pad)
def conv2d_same(image, kernel):
    kh, kw = kernel.shape
    ph, pw = kh//2, kw//2
    padded = np.pad(image, ((ph,ph),(pw,pw)), 'constant')
    H, W = image.shape
    out = np.zeros((H,W), dtype=np.float64)
    for i in range(H):
        for j in range(W):
            out[i,j] = np.sum(padded[i:i+kh,j:j+kw]*kernel)
    return out

# OpenCV (BORDER_REFLECT_101 mac dinh)
out_cv = cv2.filter2D(image, -1, kernel)
\end{lstlisting}

\begin{notebox}
\textbf{Sai khác vùng biên:} zero-padding $\to$ viền tối;\\
\texttt{cv2.filter2D} dùng reflect $\to$ tự nhiên hơn.\\
Vùng trong: kết quả \textbf{gần như giống hệt}.
\end{notebox}

% ══════════════════════════════════════════════════════════════════════════════
\section{4. Các bộ lọc làm trơn}

\subsection{Trung bình (Mean / Box Filter)}
\begin{lstlisting}[language=Python]
k_mean = np.ones((5,5)) / 25.0   # NumPy
out_cv = cv2.blur(image, (5,5))   # OpenCV
\end{lstlisting}

\subsection{Gaussian Filter}
\[
  G(x,y) = \frac{1}{2\pi\sigma^2} e^{-\frac{x^2+y^2}{2\sigma^2}}
\]
\begin{lstlisting}[language=Python]
out = cv2.GaussianBlur(image, (ksize, ksize), sigma)
# sigma=0 -> tu tinh; sigma=0.5 -> nhe; sigma=2.0 -> mo ro
\end{lstlisting}

\begin{center}\small
\begin{tabular}{@{}lll@{}}
\toprule
 & Mean & Gaussian \\
\midrule
Kernel & Đều & Giảm dần từ tâm \\
Giữ biên & Kém & Tốt hơn \\
Tăng ksize & Mờ đều, lan biên & Mờ mượt, biên mềm \\
\bottomrule
\end{tabular}
\end{center}

% ══════════════════════════════════════════════════════════════════════════════
\section{5. Bộ lọc trung vị (Median)}

\begin{lstlisting}[language=Python]
out = cv2.medianBlur(image, ksize)  # ksize phai le: 3,5,7,...
\end{lstlisting}

\begin{tipbox}
\textbf{Khi nào dùng Median?} Nhiễu Salt \& Pepper (nhiễu xung).\\
Median \textbf{robust với outlier}: giá trị 0/255 luôn nằm ở 2 đầu khi sắp xếp $\Rightarrow$ bị loại bỏ hoàn toàn. Mean/Gaussian chỉ \textbf{lan} nhiễu xung ra xung quanh.
\end{tipbox}

% ══════════════════════════════════════════════════════════════════════════════
\section{6. Khử nhiễu --- Chọn bộ lọc phù hợp}

\begin{lstlisting}[language=Python]
# Them nhieu Gaussian
def add_gaussian_noise(image, sigma=20):
    noise = np.random.normal(0, sigma, image.shape)
    return np.clip(image.astype(np.float64)+noise, 0,255).astype(np.uint8)

# Them nhieu Salt & Pepper
def add_salt_pepper(image, prob=0.02):
    out = image.copy()
    rnd = np.random.rand(*image.shape)
    out[rnd < prob/2] = 0
    out[rnd > 1-prob/2] = 255
    return out
\end{lstlisting}

\begin{center}\small
\begin{tabular}{@{}lll@{}}
\toprule
Nhiễu & Bộ lọc tốt nhất & Lý do \\
\midrule
Gaussian & GaussianBlur & Cùng phân phối \\
Salt \& Pepper & MedianBlur & Robust với outlier \\
Gaussian nhẹ & Mean filter & Đơn giản, hiệu quả \\
\bottomrule
\end{tabular}
\end{center}

% ══════════════════════════════════════════════════════════════════════════════
\section{7. Làm sắc nét (Sharpening)}

\subsection{Phương pháp 1: Laplacian}
\[
  \nabla^2 = \begin{bmatrix}0&1&0\\1&-4&1\\0&1&0\end{bmatrix}, \quad
  g = f - \lambda\,\nabla^2 f
\]
\begin{lstlisting}[language=Python]
k_lap = np.array([[0,1,0],[1,-4,1],[0,1,0]], dtype=np.float64)
lap = conv2d_same(image, k_lap)
sharp = np.clip(image.astype(np.float64) - lam*lap, 0,255).astype(np.uint8)
# lam=0.3: nhe; lam=0.7: vua; lam>=1.2 -> artifact ro
\end{lstlisting}

\subsection{Phương pháp 2: Unsharp Masking (khuyến nghị!)}
\[
  g = f + k\,\underbrace{(f - f_\text{blur})}_{\text{high-freq mask}}
\]
\begin{lstlisting}[language=Python]
blur = cv2.GaussianBlur(image, (5,5), 1.0)
hf = image.astype(np.float64) - blur.astype(np.float64)
sharp = np.clip(image.astype(np.float64) + k*hf, 0,255).astype(np.uint8)
# k=0.5: nhe; k=1.0: vua; k>=1.5 -> overshoot
\end{lstlisting}

\begin{center}\small
\begin{tabular}{@{}lll@{}}
\toprule
 & Laplacian & Unsharp Masking \\
\midrule
Nhạy nhiễu & Rất nhạy (bậc 2) & Ít nhạy (blur lọc trước) \\
Artefact & Halo mạnh & Overshoot nhẹ \\
Dễ dùng & Cần chỉnh $\lambda$ & Ít rủi ro hơn \\
\bottomrule
\end{tabular}
\end{center}

% ══════════════════════════════════════════════════════════════════════════════
\section{8. Dò biên Sobel}

\[
G_x=\begin{bmatrix}-1&0&1\\-2&0&2\\-1&0&1\end{bmatrix},\quad
G_y=\begin{bmatrix}-1&-2&-1\\0&0&0\\1&2&1\end{bmatrix}
\]
\[
  \|\nabla f\| = \sqrt{G_x^2+G_y^2}, \quad
  \theta = \arctan2(G_y, G_x)
\]

\begin{lstlisting}[language=Python]
Gx = np.array([[-1,0,1],[-2,0,2],[-1,0,1]], dtype=np.float64)
Gy = np.array([[-1,-2,-1],[0,0,0],[1,2,1]], dtype=np.float64)
Sx = conv2d_same(img, Gx);  Sy = conv2d_same(img, Gy)
mag = np.sqrt(Sx**2 + Sy**2)
# Nguong hoa
edge_map = (mag >= T).astype(np.uint8) * 255
# T=30: nhieu bien; T=60: can bang; T=90: sach but dut
\end{lstlisting}

% ══════════════════════════════════════════════════════════════════════════════
\section{9. Thuật toán Canny (5 bước)}

\begin{infobox}
\textbf{5 bước Canny:}
\begin{enumerate}
  \item \textbf{Gaussian blur} --- loại nhiễu
  \item \textbf{Sobel gradient} --- cường độ + hướng biên
  \item \textbf{NMS} --- làm mảnh biên (chỉ giữ cực đại)
  \item \textbf{Double threshold} --- phân loại strong/weak
  \item \textbf{Hysteresis} --- kết nối weak $\to$ strong
\end{enumerate}
\end{infobox}

\begin{lstlisting}[language=Python]
edges = cv2.Canny(image, threshold1=30, threshold2=80)
\end{lstlisting}

\subsection{NMS (Non-Maximum Suppression)}
\begin{lstlisting}[language=Python]
def non_maximum_suppression(G, theta):
    angle = theta * 180/np.pi
    angle[angle < 0] += 180
    Z = np.zeros_like(G)
    for i in range(1, G.shape[0]-1):
        for j in range(1, G.shape[1]-1):
            if (0<=angle[i,j]<22.5) or (157.5<=angle[i,j]<=180):
                q,r = G[i,j+1], G[i,j-1]   # 0 do
            elif 22.5<=angle[i,j]<67.5:
                q,r = G[i+1,j-1], G[i-1,j+1] # 45 do
            elif 67.5<=angle[i,j]<112.5:
                q,r = G[i+1,j], G[i-1,j]   # 90 do
            else:
                q,r = G[i-1,j-1], G[i+1,j+1] # 135 do
            if G[i,j]>=q and G[i,j]>=r:
                Z[i,j] = G[i,j]
    return Z
\end{lstlisting}

\subsection{Double Threshold + Hysteresis}
\begin{lstlisting}[language=Python]
def double_threshold(img, t_low, t_high, weak=50, strong=255):
    res = np.zeros_like(img, dtype=np.uint8)
    res[img >= t_high] = strong
    res[(img >= t_low) & (img < t_high)] = weak
    return res

def hysteresis(img, weak=50, strong=255):
    out = img.copy()
    for i in range(1, out.shape[0]-1):
        for j in range(1, out.shape[1]-1):
            if out[i,j] == weak:
                out[i,j] = strong if np.any(
                    out[i-1:i+2,j-1:j+2]==strong) else 0
    return out
\end{lstlisting}

\begin{center}\small
\begin{tabular}{@{}lll@{}}
\toprule
$(T_l, T_h)$ & Biên & Nhiễu \\
\midrule
$(20, 60)$ & Nhiều, liên tục & Nhiễu nhiều \\
$(30, 80)$ & \textbf{Cân bằng} & Vừa phải \\
$(50,120)$ & Ít, sạch & Ít, nhưng đứt \\
\bottomrule
\end{tabular}
\end{center}

\begin{tipbox}
\textbf{Vì sao hysteresis tốt hơn ngưỡng đơn?}\\
Giữ weak edge \textbf{chỉ khi} nó kết nối với strong edge $\Rightarrow$ loại nhiễu cô lập, biên liên tục hơn.
\end{tipbox}

% ══════════════════════════════════════════════════════════════════════════════
\section{10. So sánh Spatial vs Frequency Filtering}

\begin{center}\small
\begin{tabular}{@{}lll@{}}
\toprule
 & Spatial & Frequency \\
\midrule
Thao tác & Kernel trượt & Nhân phổ \\
Gaussian blur & GaussianBlur & Gaussian LPF \\
Giống nhau? & Không hoàn toàn & (ideal $\neq$ Gaussian) \\
Trực quan & Dễ & Cần biết FFT \\
Ứng dụng & Thông dụng & Lọc chọn lọc \\
\bottomrule
\end{tabular}
\end{center}

\begin{notebox}
Gaussian spatial blur $\neq$ Circular LPF vì:\\
Circular LPF lý tưởng cắt ngưỡng đột ngột $\Rightarrow$ ringing.
Gaussian blur $\leftrightarrow$ Gaussian LPF (mượt, không ringing).
\end{notebox}

% ══════════════════════════════════════════════════════════════════════════════
\section{11. Mini-Project: Pipeline hoàn chỉnh}

\begin{center}
\textbf{Ảnh gốc $\to$ Nhiễu $\to$ Khử nhiễu $\to$ Sắc nét $\to$ Biên}
\end{center}

\begin{lstlisting}[language=Python]
img_clean = img_edge.copy()   # anh sach ban dau

# === Pipeline 1: Gaussian noise ===
noisy_g = add_gaussian_noise(img_clean, sigma=25)
denoised_g = cv2.GaussianBlur(noisy_g, (5,5), 1.0)
blur_tmp = cv2.GaussianBlur(denoised_g, (5,5), 1.0)
hf = denoised_g.astype(np.float64) - blur_tmp.astype(np.float64)
sharp_g = np.clip(denoised_g.astype(np.float64)+1.0*hf, 0,255).astype(np.uint8)
edges_g = cv2.Canny(sharp_g, 30, 80)

# === Pipeline 2: Salt & Pepper noise ===
noisy_sp = add_salt_pepper(img_clean, prob=0.03)
denoised_sp = cv2.medianBlur(noisy_sp, 5)
blur_tmp2 = cv2.GaussianBlur(denoised_sp, (5,5), 1.0)
hf2 = denoised_sp.astype(np.float64) - blur_tmp2.astype(np.float64)
sharp_sp = np.clip(denoised_sp.astype(np.float64)+1.0*hf2, 0,255).astype(np.uint8)
edges_sp = cv2.Canny(sharp_sp, 30, 80)
\end{lstlisting}

\begin{center}\small
\begin{tabular}{@{}lllll@{}}
\toprule
Nhiễu & Khử & Sắc nét & Biên \\
\midrule
Gaussian $\sigma$=25 & GaussianBlur & Unsharp $k$=1 & Canny(30,80) \\
S\&P $p$=0.03 & MedianBlur & Unsharp $k$=1 & Canny(30,80) \\
\bottomrule
\end{tabular}
\end{center}

\begin{notebox}
\textbf{Trade-off 3 chiều --- Quy tắc vàng:}
\begin{itemize}
  \item Khử nhiễu mạnh $\Rightarrow$ ít nhiễu nhưng \textbf{mất biên yếu}
  \item Sắc nét mạnh $\Rightarrow$ rõ biên nhưng \textbf{khuếch đại nhiễu}
  \item Ngưỡng Canny cao $\Rightarrow$ sạch nhưng \textbf{biên đứt đoạn}
\end{itemize}
$\Rightarrow$ \textbf{Không có tham số ``tốt nhất'' chung}, phải cân bằng cả 3 cho từng ảnh!
\end{notebox}

% ══════════════════════════════════════════════════════════════════════════════
\section{12. Bảng hàm OpenCV nhanh}

\begin{center}\scriptsize
\begin{tabular}{@{}ll@{}}
\toprule
Hàm & Mục đích \\
\midrule
\texttt{cv2.blur(img,(k,k))} & Lọc trung bình \\
\texttt{cv2.GaussianBlur(img,(k,k),s)} & Lọc Gaussian \\
\texttt{cv2.medianBlur(img,k)} & Lọc trung vị \\
\texttt{cv2.filter2D(img,-1,kernel)} & Tích chập tổng quát \\
\texttt{cv2.Canny(img,t1,t2)} & Dò biên Canny \\
\texttt{np.fft.fft2(img)} & FFT 2D (complex) \\
\texttt{np.fft.fftshift(F)} & DC ra tâm \\
\texttt{np.fft.ifft2(F)} & IFFT 2D \\
\texttt{np.fft.ifftshift(F)} & Đảo fftshift \\
\bottomrule
\end{tabular}
\end{center}

% ══════════════════════════════════════════════════════════════════════════════
\section{13. Công thức cần nhớ \& Tips}

\begin{center}\scriptsize
\begin{tabular}{@{}ll@{}}
\toprule
Tên & Công thức \\
\midrule
DFT & $F(u,v) = \sum_x\sum_y f\,e^{-j2\pi(ux/M+vy/N)}$ \\
Tích chập & $(f*h)(x,y)=\sum_i\sum_j f(x{-}i,y{-}j)h(i,j)$ \\
Conv. theorem & $f*h \Leftrightarrow F\cdot H$ \\
Gaussian & $G=\frac{1}{2\pi\sigma^2}e^{-(x^2+y^2)/2\sigma^2}$ \\
Gradient & $\|\nabla f\|=\sqrt{G_x^2+G_y^2}$ \\
Unsharp & $g=f+k(f-f_\text{blur})$ \\
Laplacian & $g=f-\lambda\nabla^2 f$ \\
\bottomrule
\end{tabular}
\end{center}

\begin{tipbox}
\textbf{7 tips thực hành:}
\begin{enumerate}
  \item Luôn chuyển \texttt{uint8 $\to$ float64} trước khi tính
  \item \texttt{np.clip(result,0,255).astype(np.uint8)} tránh tràn số
  \item Phổ FFT dùng \texttt{log(1+|F|)}, không để DC nuốt hết
  \item Zero-pad $\to$ viền tối; reflect-pad $\to$ tự nhiên hơn
  \item So sánh NumPy vs OpenCV: \texttt{np.abs(a-b)} diff map
  \item \textbf{Khử nhiễu TRƯỚC khi dò biên} --- không bao giờ ngược
  \item \textbf{Median cho S\&P, Gaussian cho Gaussian noise} --- quy tắc vàng
\end{enumerate}
\end{tipbox}

\end{multicols}
\end{document}
